% Standalone section module; requires amsmath, amssymb, amsthm and the
% theorem/proposition/corollary/remark environments. Bibliography: HK2017,HKO2026.
\section{First variation under simultaneous facet splittings}
\label{sec:multicut-variation}

Fixed-facet variations do not describe the birth of several facets at once.
Here the relevant first-order datum is a convex function on each old facet:
it records the maximum depth of the competing cuts. We obtain its exact
capacity response from labeled occupation measures of shortest polygonal
characteristics. The resulting convex-order criterion is a criterion for
\emph{first-order stationarity}. It does not prove a Hausdorff neighborhood
of local maximality for the pentagon counterexample.

Let $K=\{x\in\mathbb R^{2n}:a_i\cdot x\leq1,\ 1\leq i\leq m\}$
be a bounded polytope containing the origin in its interior, with genuine
facets $F_i$. Use $\omega(u,v)=\langle Ju,v\rangle$, where $J^2=-I$.
For a finite word of facet labels and weights, put
\begin{equation}\label{eq:mc-word}
 Q(\sigma,\beta)=\sum_{r<s}\beta_r\beta_s
 \omega(a_{\sigma_r},a_{\sigma_s}),\qquad
 \sum_r\beta_ra_{\sigma_r}=0,\quad
 \sum_r\beta_r=1,\quad\beta_r\geq0.
\end{equation}
Haim--Kislev's formula~\cite{HK2017} gives
$Q_*:=\max Q=1/(2c_{\mathrm{EHZ}}(K))$ using words without repeated labels.
Repeated words also satisfy $Q\leq Q_*$: their unit-period polar loop has
velocity $2Ja_{\sigma_r}$ for time $\beta_r$, action $2Q$, and dual energy
$\frac14\int h_K(-J\dot z)^2=1$. Normalizing its action to one in the dual
action principle gives $c_{\mathrm{EHZ}}(K)\leq1/(2Q)$. Translating the loop
to zero mean changes neither quantity. For redundant rows the support values
are at most one, so the same inequality holds.

We impose the following condition:
\begin{quote}
Every finite maximizing word, after deleting its zero weights, has
$\operatorname{rank}\{(a_{\sigma_r},1)\}_r=2n+1$.
\end{quote}
We call this the \emph{full-rank support condition}. It concerns the constraint
matrix, not the Hessian of the quadratic program.

A maximizing word reconstructs a shortest boundary characteristic, as shown
below. Parametrize it on $[0,1]$ with velocity $2TJ a_{\sigma_r}$ for time
$\beta_r$, where $T=c_{\mathrm{EHZ}}(K)$. Its labeled occupation measure
$\mu_\gamma$ is a probability on the disjoint union
$X=\bigsqcup_iF_i$: time on a segment with label $i$ is assigned to that
copy of $F_i$, including at ridges. Let $\mathcal M$ be the weak closure of
the measures of \emph{all finite maximizing words}, allowing repetitions.
The normalized cone-volume probability is
\begin{equation}\label{eq:mc-nu}
 d\nu|_{F_i}=\frac{dS}{2n\,\operatorname{Vol}_{2n}(K)|a_i|}.
\end{equation}

\begin{theorem}[Simultaneous-splitting derivative]\label{thm:mc-derivative}
Assume the full-rank support condition. For each facet choose a nonempty
finite list $D_{ij}\in\mathbb R^{2n}$ and set
\[
 K_t=\{x:(a_i+tD_{ij})\cdot x\leq1\text{ for every }i,j\},\qquad
 f_i(x)=\max_jD_{ij}\cdot x\quad(x\in F_i).
\]
For this fixed finite list, the right derivative exists and equals
\begin{equation}\label{eq:mc-derivative}
 \left.\frac{d}{dt}\right|_{0+}
 \log\frac{c_{\mathrm{EHZ}}(K_t)^n}{n!\operatorname{Vol}_{2n}(K_t)}
 =2n\left(\int_X f\,d\nu-
                  \max_{\mu\in\mathcal M}\int_X f\,d\mu\right).
\end{equation}
In particular, the hypothesis holds for every Lagrangian product of two
planar polygons having no antipodal facet normals, including the HKO product.
\end{theorem}

\begin{proof}
We give the sensitivity and placement arguments explicitly, since neither
requires a smooth optimizing branch.

\emph{A local chart at every limiting optimizer.}
For the fixed augmented row list, the moving optimum $q(t)$ is the maximum
over finitely many permutations and their feasible weights. This remains
true with redundant or repeated presentation rows: genuine facets give one
inequality and the polar-loop bound above gives the other. Along any sequence
$t\downarrow0$, compactness of the weight simplex and finiteness of the words
give maximizing weights $\beta_t\to\beta_0$ in one word. Capacity continuity
makes $\beta_0$ a base maximizer.
Choose $2n+1$ positive coordinates $I$ forming an invertible constraint minor,
and denote the remaining weights by $z$. With columns
$B_r(t)=(a_{\sigma_r}+tD_r,1)$, solve
\[
 \beta_I(t,z)=B_I(t)^{-1}\bigl(e-B_J(t)z\bigr),
 \qquad e=(0,\ldots,0,1)^T.
\]
The solved weights stay positive near $(0,z_0)$; the other coordinates merely
satisfy $z\geq0$. In particular $\beta(0,z_t)$ is base-feasible.
Writing $F(t,z)=Q_t(\beta(t,z))$, we obtain
\[
 q(t)-Q_*\leq F(t,z_t)-F(0,z_t)
             =t\,\partial_tF(0,z_0)+o(t).
\]
Conversely the section with $z=z_0$ fixed gives the reverse lower bound for
every base optimizer. Thus $q'(0+)$ exists and is the maximum of these
section derivatives. This comparison also covers singular Hessians and
continua of maximizing weights.

\emph{Placement of a maximizing word.}
Delete zero weights and write $a_r=a_{\sigma_r}$,
$s_r=\sum_{k<r}\beta_ka_k$, and $p_r=\partial_{\beta_r}Q$.
Stationarity and full rank give unique multipliers satisfying
\[
 p_r=\lambda\cdot a_r+\eta,
 \qquad \eta=2Q_*=1/T,
 \qquad p_r=-2\omega(a_r,s_r)=2a_r\cdot Js_r.
\]
The vertices $x_r=-T\lambda+2TJs_r$ therefore satisfy $a_r\cdot x_r=1$.
The segment to $x_{r+1}$ has direction $Ja_r$, so its active facet equation
holds throughout. To obtain every other facet inequality, insert an arbitrary
normal $a_k$ at any point of the loop, subdividing its current segment if
necessary. Give the inserted row weight $\varepsilon\geq0$ and compensate
with the positive minor. Its reduced objective derivative is
\[
 p_{\rm new}-\lambda\cdot a_k-\eta
   =\frac{a_k\cdot x-1}{T}.
\]
Global optimality among repeated words makes this nonpositive. Hence the
whole loop lies in $K$, on its labeled facets, and has action $T$.

\emph{The midpoint formula.}
Differentiate a feasible section of this word with row velocities $D_r$.
The differentiated constraints give
$\sum_ra_r\beta'_r=-\sum_r\beta_rD_r$ and $\sum_r\beta'_r=0$.
The direct row derivative of $Q$ is
\[
 \sum_r\bigl[-2\beta_r\omega(D_r,s_r)
              -\beta_r^2\omega(D_r,a_r)\bigr].
\]
Adding the weight derivative $-\sum_r\beta_r\lambda\cdot D_r$ yields
\begin{equation}\label{eq:mc-midpoint}
 Q'=\frac1T\sum_r\beta_rD_r\cdot x_r^{\rm mid},
 \qquad x_r^{\rm mid}=x_r+T\beta_rJa_r.
\end{equation}
The self term in the preceding display is essential: the evaluation point
is the segment midpoint, not its starting point.
For every optimizer in the finite bank, $D_r\cdot x\leq f_i(x)$ on its
segment labeled $i$. Affine midpoint integration therefore gives
$q'(0+)\leq T^{-1}\max_{\mathcal M}\int f\,d\mu$.
For the reverse inequality, take any finite maximizing base word and
subdivide every segment at the finitely many switches of the maximum
$f_i=\max_jD_{ij}\cdot x$. Assign a maximizing row on each resulting interval.
After deleting zero-length pieces, this is a positive repeated word with the
same base objective and full rank. A feasible continuation exists, and the
polar-loop inequality bounds its objective by $q(t)$.
Equation~\eqref{eq:mc-midpoint} now integrates $f$ exactly. Taking the supremum
and then weak closure gives
\[
 q'(0+)=2Q_*\max_{\mu\in\mathcal M}\int f\,d\mu,
 \qquad \frac{c'(0+)}{c(K)}=-2\max_{\mu\in\mathcal M}\int f\,d\mu.
\]
The maximum over $\mathcal M$ exists because $X$ is compact.

\emph{Volume and the product hypothesis.}
Away from the ridges, the radial gauge of $K_t$ has derivative $f_i$ on
$F_i$. The radial volume formula and dominated convergence, using the fixed
finite row list, give
\[
 V'(0+)=-\sum_i\int_{F_i}\frac{f_i}{|a_i|}\,dS
        =-2nV\int f\,d\nu.
\]
This proves~\eqref{eq:mc-derivative}.
For a planar Lagrangian product, a positive-$Q$ word must use both factors.
Closure separates into the two planes. A positive closed combination of
nonzero planar normals without antipodal pairs spans its plane; hence the
active normals span $\mathbb R^4$. If $v\cdot a_r+s=0$ on their positive
support, weighting and summing gives $s=0$, then $v=0$. The augmented rank is
therefore five, also for repeated words.
\end{proof}

\begin{corollary}[A common convex-order ensemble]\label{cor:mc-convex-order}
Under the theorem's hypothesis, the following are equivalent:
\begin{enumerate}[label=\textup{(\roman*)}]
\item Every fixed finite simultaneous-splitting direction has nonpositive
first systolic variation.
\item The weakly closed convex hull of $\mathcal M$ contains a probability
$\bar\mu$ such that
$\int f\,d\bar\mu\geq\int f\,d\nu$ for every continuous function convex
on each facet copy.
\end{enumerate}
It is enough in \textup{(i)} to use inward cuts retaining every old facet.
Any such $\bar\mu$ has the same mass and first moment as $\nu$ on each facet.
\end{corollary}
\begin{proof}
Finite maxima of affine functions are uniformly dense in the continuous
convex functions on a compact facet. Every affine function there is $D\cdot x$
because $a_i\cdot x=1$. Thus the theorem turns (i) into
$\sup_{\mathcal M}\int f\,d(\mu-\nu)\geq0$ for every facetwise convex $f$.
Let $B=\overline{\operatorname{conv}}\mathcal M$. For a finite list
$f_1,\ldots,f_k$, the compact convex set of vectors
$(\int f_j\,d(\mu-\nu))_{j=1}^k$, $\mu\in B$, meets the positive orthant.
Otherwise separation supplies nonnegative coefficients $\alpha_j$ whose
sum $\sum_j\alpha_jf_j$ contradicts the preceding inequality.
Compactness and the finite-intersection property now give one common
$\bar\mu$ for all $f$. The converse is immediate.
Adding one common large constant to all facet profiles makes them
nonnegative, allowing the zero row to be retained; the difference of two
probability integrals annihilates that constant. Finally test both signs of
facetwise affine functions to obtain the mass and first-moment identities.
\end{proof}

The common-ensemble requirement is stronger than finding a different
successful orbit for each single cut. Nor does first-order nonincrease give
an unrestricted neighborhood theorem: zero derivatives and estimates whose
constants deteriorate with cut count or slope still require analysis.

\subsection{Two limitations of finite orbit reductions}

The HKO body is $H=P\times_LR(-\pi/2)P$, where $P$ is the circumradius-one
regular pentagon. Put $\alpha=\pi/5$, $\tau=\tan\alpha$, and
\[
 a_0=(1,\tau,0,0),\qquad w=(-\sin\alpha,\cos\alpha).
\]
On $F_0=E_0\times R(-\pi/2)P$, the functional $\ell(q,p)=w\cdot p$
attains its maximum one exactly on $E_0\times\{w\}$.

\begin{proposition}[No finite pure-word ensemble is universal]
No finite convex combination of occupation measures of finite pure
facet-word characteristics of $H$ dominates $\nu$ for every facetwise convex
profile. Here a pure segment has velocity a positive multiple of one
$Ja_i$ and carries that label; shortestness is not required.
\end{proposition}
\begin{proof}
On a segment labeled zero, $\ell$ is constant since
$w\cdot(1,\tau)=0$. Its value is strictly below one: a nonconstant
p-segment cannot lie in the singleton exposing slice at $w$.
A finite bank therefore has a common bound $\ell\leq h<1$ on all its
zero-labeled segments. For $h<r<1$, the profile
$f_0=(\ell-r)_+$, $f_i=0$ for $i\ne0$, vanishes on the bank but has
positive $\nu$-integral. It is realized by the single extra row
$a_0+t(-r,-r\tau,w_1,w_2)$.
More explicitly, for $\delta=1-r<1-\cos(2\alpha)$, the pentagon cap
has slice width $2z\cot\alpha$, so
\[
 \int_Xf\,d\nu=\frac{\cot\alpha}{30\operatorname{area}(P)}\delta^3>0.
\]
\end{proof}

The existing at-most-eleven-facet HKO proof already identifies the failure
of fixed positive-weight banks at redundant limits and uses cut-dependent
families with vanishing weights. The proposition gives a geometric form of
that obstruction; it does not contradict that theorem. Infinite ensembles
and weak limits remain available in Corollary~\ref{cor:mc-convex-order}.
For an action-$T$ word, the wedge at $w$ also gives the quantitative bound
$1-\ell\geq T\beta\sec\alpha\tan\alpha$ on a zero-labeled segment of
weight $\beta$, by comparing its length with the wedge's slice width.

A second obstruction concerns eliminating repeated normals by moving their
weights to an endpoint. That operation need not maximize a convex profile's
occupation response, even along a family of genuine HKO minimizers.
For the following explicit example set $r=\sqrt5$, $\tau=\sqrt{5-2r}$,
$u=(3-r)/2$, $v=\tau(1+r)/2$, and $s=r-1$. The ten normals are
\[
\begin{split}
 &(1,\tau,0,0),\ (-u,v,0,0),\ (-s,0,0,0),\ (-u,-v,0,0),\ (1,-\tau,0,0),\\
 &(0,0,\tau,-1),\ (0,0,v,u),\ (0,0,0,s),\ (0,0,-v,u),\ (0,0,-\tau,-1).
\end{split}
\]

\begin{proposition}[An interior response maximum]\label{prop:mc-interior}
Let $b=(5-r)/20$ and consider
\[
 \sigma=(0,6,7,3,2,8,3,9),\qquad 0\leq q\leq1,
\]\[
 \beta(q)=\left(\frac r{10},\frac{5+r}{40},\frac{5-r}{40},
 b(1-q),b,\frac{5-r}{40},bq,\frac{5+r}{40}\right).
\]
Every member is a maximizing word. For the facetwise convex profile
$f_3(x)=|x-c_3|^2$, $f_i=0$ otherwise, where
$c_3=(-1/4,\tau(\tau^2-9)/8,0,0)$ is the facet centroid, its response is
\[
 g(q)=\int f\,d\mu_q
 =\frac{2-r}{4}q^2+\frac{-19+9r}{16}q+\frac{195-77r}{240}.
\]
Its maximum occurs at $q_*=(7-r)/8$ and satisfies
$g(q_*)-\max\{g(0),g(1)\}=(r-1)/128>0$.
\end{proposition}
\begin{proof}
Substitution gives closure and $Q=\tau/5$ at both endpoints. The objective
is affine in the transfer between two identical normals, since
$\omega(a_3,a_3)=0$, so it is constant throughout. This is the HKO optimum
by~\cite{HKO2026}; the placement argument above applies.
For a direct calculation of the response, the two row-three segment
midpoints are
\[
 m_1(q)=\frac14\bigl(-1-r,-\tau(1+r),\tau(r+q),1+rq\bigr),\qquad
 m_2(q)=\frac14\bigl(-1,-\tau(2+r),\tau(q-3-r),1-r+rq\bigr).
\]
Their weights are $b(1-q)$ and $bq$. Exact integration of a quadratic along
a straight segment of weight $\beta$ gives
$\beta(|m-c_3|^2+T^2\beta^2|a_3|^2/3)$.
Here $T=5/(2\tau)$ and $T^2b^2|a_3|^2/3=(5-r)/24$.
Inserting the displayed midpoints yields $g$; differentiating this
strictly concave quadratic proves the claimed maximum and gap.
\end{proof}

Uniform approximation of this profile by finite maxima of affine functions
preserves the strict endpoint gap. Thus the phenomenon also occurs for a
finite inward split bank. It refutes the proposed convexity argument for
deleting repetitions; it does not show that the two displayed endpoint words
are best among \emph{all} simple HKO words, and does not exhibit an improving
systolic variation.
